\section{Dialect design}

The \texttt{npu} dialect is the tensor level. Values are builtin ranked fp32
tensors, which keeps interoperability with the MLIR infrastructure so that generic
folding and dead code elimination apply without op specific support. Convolution
and matmul take an optional bias operand and carry a fused activation attribute
rather than exploding into a family of named fused ops. This mirrors the fused
operator style of TFLite and gives the fusion pass a single clean target. Every op
that computes a pure function of its operands carries the \texttt{Pure} trait, and
each op with shape or type constraints has an ODS verifier. The dialect reference
is generated from the op definitions so it cannot drift.

The \texttt{npuisa} dialect is the instruction level, after memory placement. A
value of \texttt{!npuisa.buffer} type is a scratchpad resident tensor. DRAM
tensors move to and from the scratchpad through \texttt{dma\_load} and
\texttt{dma\_store}, and compute instructions read and write buffers only. The
instruction stream is straight line, since inference graphs are static DAGs, so
the instruction set has no branches. Each buffer producing instruction carries an
address attribute that the allocator fills in.
